那已经是六七年前的事情了，那时我还在北京工作，周五傍晚，外面窸窸窣窣下起了小雨，那时的四号线还没有现在这么挤，但周五的晚高峰照例是没有座位的。刚走进地铁站，就感觉左肩膀被人拍了一下。

“亮，没想到会在这里遇到你。” 一阵温柔的女声传来，我转过头，眼前这人穿一件贴身的短上衣和很宽的裤子，头发松松地垂在肩上，正朝我笑。我怔了一下，才认出来那是理佳。明明七年没有见过，可是理佳的笑容却一点不显得陌生。

理佳是我的高中同学，高中毕业后我就再也没见过她，当时她坐在后排，而我坐在前排，我们彼此当然认识，但也仅限于认识，若不是她突然喊住我，我恐怕还难以将她和坐在后排靠窗的那个女孩对应起来。

“嗨，理佳，真是没想到我们能在这里见面。我们从高中毕业后就从未见过吧。”

“是啊，高中毕业后就再也没有见过，到现在都七年时间了呢。” 理佳笑吟吟地望向我，我突然有一阵恍惚，高中时我习惯了一个人，她周围却永远围着一群人，我们之间隔着的，不只是一个教室的距离，若非此次突然遇到，恐怕之后便再也没有交集了。

“你也坐四号线吗？”

“对，我也坐四号线，我要坐到国家图书馆换乘，你坐到哪里，我们顺路吗？”

“我和你一个方向，这样我们路上可以聊聊天。”

我们寒暄了一阵，地铁就来了，不出意料的是地铁上没有座位，晚高峰的地铁上照例没有座位，好在这趟地铁不算拥挤，我们不至于在人群的呼吸声中进行肉搏。理佳在地铁中一处人少的角落中站定，我随即走了过去。

“你现在在做什么呢？是工作了，还是还在读书？” 理佳问道。

“我现在在一家公司做电脑工程师，你呢，你现在在做什么？”

“我现在在一所高中做教师。”

“你现在怎么做老师了，我记得你高中毕业时候选的专业不是经济金融这一类吗，怎么现在突然当起了老师？”

“唉，高中时候还是年轻，不知道想要什么，家里面觉得金融的就业很好，当时的分数也足够，就选了这个专业，可是本科时候去过投行，也去过大厂实习，感觉工作压力太大，自己也实在不喜欢，所以后来就去考了教师资格证当老师了，感觉还是有更多休闲的时间更适合我。”

“那你觉得当老师怎么样，你喜欢这份工作吗？” 我问道。

“也谈不上喜欢或者不喜欢，感觉工作也都差不多，其实当时就是因为大家觉得这份工作安稳，我就选了。平时备课和讲课倒是还好，跟上学的时候做题也都差不多。最让我犯难的时候是开家长会，任课教师要挨个跟家长沟通，天哪，你知道吗，那个时候我总觉得我是在装大人，我根本不知道该用什么立场去教导他们，就只能顺着话说，机械地夸奖，安慰。还要面对他们奇奇怪怪的焦虑，甚至有人跑来问我，能不能让孩子用 AI 帮忙写作文......我觉得我快被掏空了。”

“或许这才是工作的真相吧，”，我苦笑一声，“即使是从前非常迷恋的兴趣，当将它变为工作后还能保留几分呢？”

“啊... 这样啊，”，理佳眼中闪过一丝同情，不过她的同情究竟是真的理解，还是仅仅是成年人的得体的社交辞令，我却并不了解。

“那真是可惜呢，”，她用一种轻柔的声音安慰我，“不过，既然你找到了一份稳定又不错的工作，等熬过了最初阶段，周末还是可以把原来的兴趣爱好捡起来啊。”

她说到这里，偏了偏头，“不过说起来...... 我记得你高中时候不是在写小说吗？我记得当时你还说过要给什么严肃文学杂志投稿，你现在...... 还在写吗？”

我愣了一下，话说，她是怎么知道我在写小说的事情呢？我记得我在高中时候写的小说很少给人看，或许是柏宇给她讲的吧。

“早就不写了，自打高中之后就不写了，偶尔写一些分析类的文章，小说一直没什么灵感，上班后也没什么时间。” 我的声音比想象中的还要短促。

“那很可惜啊，我听柏宇说你写小说也很好，他很喜欢读你的小说呢。”

柏宇是我高中时期为数不多的朋友，而且他也是我们中最聪明的一个，他高中时候成绩很好，读上了名校，就当我们以为他会大展前程的时候，他会选择回到我们的家乡城市工作 ...... 果然是他。

她顿了一下，像突然想起了什么，“话说你现在还和柏宇有联系吗，他现在怎么样了？”

“我们有时会见一下，但次数也不太多，我们过年吃了顿饭。虽然大家都觉得他回家乡工作很可惜，但其实和他聊起来感觉他还蛮开心的。”

“蛮开心的... 原来你们是这样认为的吗？” 理佳回复道，若有所思。

聊着聊着我们到了换乘站，我和她说我要去一趟卫生间，让她稍微等一下。“把你的伞给我拿着吧”，她说。

她接过伞，手指碰到我的，那触感比我想象的柔软，那一瞬间，地铁的雨气仿佛变了质地，不再是阴冷的潮湿，而是某种更温热的雾气。她低头把伞收好，头发从肩上滑下来。

出来之后，我发现她并不在门口。我四处张望，转头却发现她靠在墙上，左手拿着我的伞放在腿上，右手掐着一根烟。这使得我对她的印象有些陌生，隔着升腾的烟雾，原本那个坐在后排靠窗的剪影，在我心里突然破碎。她微微蹙着眉，面容疲惫。

但我心底却没有升起被欺骗的愤怒。相反，我看着她夹着烟的微颤的手指，一种极其隐秘的、自负的怜悯涌了上来。我觉得我窥见了她真实伤口的一角。我甚至有些迷恋她此刻这种陌生而危险的形象。

她看到我，将烟在垃圾桶上的烟灰缸里捻灭。“走吧，”她说，又恢复了方才的笑容，把伞递还给我。

换乘之后，我站在车厢里，心跳还是很重，颇有一种通宵喝咖啡的感觉，隧道中的灯一条一条从窗外过去，我恍惚间又看到她的头发从肩上滑下来。我的呼吸逐渐变得粗重，我忽然想，如果我现在坐回去，去她的换乘站，或许她还在那里？
